\documentclass[11pt]{article}
\usepackage[margin=0.9in]{geometry}
\usepackage{amsmath,amssymb,amsthm,booktabs,mathrsfs}
\usepackage[T1]{fontenc}
\usepackage{lmodern,microtype}
\usepackage[colorlinks=true,linkcolor=blue,urlcolor=blue]{hyperref}
\newtheorem{theorem}{Theorem}[section]
\newtheorem{lemma}[theorem]{Lemma}
\newtheorem{proposition}[theorem]{Proposition}
\newtheorem{corollary}[theorem]{Corollary}
\theoremstyle{remark}
\newtheorem{remark}[theorem]{Remark}
\newcommand{\norm}[1]{\left\|#1\right\|}
\title{Toward a four-dimensional Yang--Mills construction\\
\large Notebook 21: local classical actions through arbitrary refinement}
\author{Research working document}
\date{17 October 2026}
\setlength{\emergencystretch}{5em}
\begin{document}
\maketitle
\begin{abstract}
We localize the actual homogeneous constrained classical action
uniformly through the preceding refinement $M=L^k$. Smooth
physical cells and the normalized fine pairing give uniform
mixed-insertion bounds for the concrete Wilson and corrected
constraint vertices. A contraction in a rooted support algebra
then constructs the minimizer and multiplier as convergent
families of functions of specified background cells. The
action-envelope identity gives strictly supported scalar
coefficients; their Hessians are individually symmetric local
precision pieces. The algebra admits an arbitrary fixed
support-cardinality weight, at the cost of a smaller uniform
background ball. Substitution into the exact fixed-$L$ source
chart gives the localized classical remainder and physical
precision needed by the overlap integration theorem. This
extends the bare small-field integral of Notebook 20 to one
integration of the classical reference action at any earlier
classical depth. Generated quantum interactions and large-field
histories still require an induction theorem.
\end{abstract}

The distinction between the two kinds of depth is explicit:
$S_M$ is obtained by classical minimization, while a full
quantum effective action also contains the corrections from
every earlier fluctuation integral. The theorem here controls
the former and one subsequent small-field integration. All
strict support claims concern the original independent
current-background variables.


% Notebook 21 concrete module. No preamble. Prefix cv-.
\section{Concrete classical vertices uniform through refinement}
\label{cv-section}

The fine spacing is $\eta=M^{-1}$, with $M=L^j$, fixed odd
$L$, dimension four and current spacing one. Fine pairings
use $\eta^4$ times counting measure; current pairings use
counting measure. All constants below are uniform in $M$ and
compatible periodic volume. The conclusion is a set of concrete
inputs for a rooted support expansion, not an assumption that
one weighted response estimate already gives a sum over all
input insertions.

\subsection{Complete physical field spaces and known bounds}
\label{cv-spaces-section}

Fix $0<\alpha<1$ and use complete fine-field spaces
\begin{equation}\label{cv-spaces}
 \begin{split}
 X_\eta=\mathcal N_\alpha:\quad
 \norm A_X&=\norm A_\infty+\norm{\nabla_\eta A}_\infty
                         +[\nabla_\eta A]_\alpha,\\
 Y_\eta&=C^\alpha\text{ of fine jet coefficients},\\
 Z_\eta&=\ell^\infty\text{ of fine coefficients},\qquad
 E=\ell^\infty\text{ of current bond fields}.
 \end{split}
\end{equation}
Local H\"older norms use physical distances and fixed cell
collars. Color and orientation spaces have fixed finite
dimension. The plaquette jet $S=(I,\nabla_\eta)$, with the
needed finite component selections, maps $X_\eta$ to $Y_\eta$
uniformly. No sharp cell indicator is differentiated.

The proved Wilson-force module gives
\[
 \mathcal S_\eta(A)=\eta^4\sum_p\mathcal L_\eta((SA)_p),
 \qquad N(A)=S^Tf_\eta(SA),
\]
where $\mathcal L_\eta$ is uniformly analytic in the jet,
its quadratic part is the Maxwell density, and $f_\eta$ starts
at degree two. The full corrected constraint of Notebook13 is
\[
 \mathcal F(A)=QA+\mathcal R(A),\qquad Q=Q_M^{\rm st}.
\]
It is analytic on a refinement-uniform supremum ball and has
a fixed physical dependency radius. Notebook15 also proves
uniform holomorphy of $D\mathcal R(A)^T$ in the physical
bounded-row operator norm, with norm $O(\norm A_\infty)$.
The transpose includes the physical fine measure.

The weighted complete-field bounds for $H,\mathcal C$ and
their differentiated compositions are recorded explicitly in
Section~\ref{cv-linear-section}. They include the
equal-exponent constrained Schauder theorem of Notebook15.
The scalar action of the ordinary homogeneous minimizer is
exactly the action at the corrected fixed-slice minimizer by
Notebook15's gauge transfer. This does not identify the two
constraint derivatives.

The underlying corrected one-step coordinate map is specified
at Raw I, lines 24320--24325 and 24344--24369, with the
ordinary/corrected bridge at 47076--47087. Its physical
scaling is $(Lh)^{-1}F_L(hA)$. The Wilson density and pairing
use the conventions already matched in Notebooks14--15.
No additional claim about the full original orbit domains is
made here.

\subsection{Smooth cell spaces without an extension assumption}
\label{cv-cells-section}

Choose a smooth partition of unity $\{\chi_i\}$ indexed by
unit current cells, with support in fixed enlarged cells
$S_i$, uniformly bounded derivatives through order three, and
uniformly finite overlap. For example, divide a nonnegative
compactly supported smooth bump by the sum of its integer
translates. Periodize on compatible tori, counting repeated
supports with multiplicity.

For each fine field type let $E_i$ be the Banach subspace of
\emph{complete} fields supported in $S_i$, with its global
$X_\eta$, $Y_\eta$ or $Z_\eta$ norm. Write $I_i$ for inclusion
and $P_iu=\chi_i u$. For current fields use coordinate
restriction and insertion in one current cell. Then
\begin{equation}\label{cv-partition}
 u=\sum_iI_iP_iu,\qquad
 \norm{P_iu}_{E_i}\le C\norm u_{\text{fixed collar of }i}.
\end{equation}
The reconstruction
\[
 \mathscr R((u_i))=\sum_iI_iu_i
\]
is bounded from the block supremum space to the complete
global space, uniformly in $\eta$ and volume. In particular,
$\mathscr RP=I$ and $P\mathscr R$ is a bounded projection.
Arbitrary restrictions on a cell are not extended by an
unproved extension theorem; the spaces $E_i$ already consist
of supported complete fields.

For the only delicate norm $X_\eta$, the exact product rule is
\begin{equation}\label{cv-product-rule}
 \nabla_\eta(\chi A)(x)
 =\chi(x+\eta e)\nabla_\eta A(x)
          +(\nabla_\eta\chi)(x)A(x).
\end{equation}
The sampled smooth functions and their first differences have
uniform physical $C^\alpha$ norms. Telescoping controls the
local $C^\alpha$ norm of $A$ by $\norm A_X$ without an inverse
mesh factor. The H\"older product inequality proves
\eqref{cv-partition}. Finite overlap, also for pairs of points
in fixed collars, proves the reconstruction bound. The
arguments for $Y_\eta$ and $Z_\eta$ are simpler.

\subsection{A summed mixed-insertion vertex bound}
\label{cv-vertex-section}

For a finite set of current cell labels $S$, let $\tau(S)$ be
the minimum length of a connected current-lattice graph
containing $S$, allowing unmarked auxiliary vertices. This is
graph Steiner length. It is not the minimum-spanning-tree
length on the terminals alone.

For a local analytic map $F$, expand after reconstruction:
\begin{equation}\label{cv-vertex-expansion}
 P_iF(\mathscr Ru)=
 \sum_{n\ge0}\sum_{j_1,\ldots,j_n}
 V_{n;i;j_1,\ldots,j_n}[u_{j_1},\ldots,u_{j_n}].
\end{equation}
The coefficients contain $1/n!$. A fixed finite direct sum
handles several field types. Define
\begin{equation}\label{cv-vertex-norm}
 v_n(\kappa)=\sup_i\sum_{j_1,\ldots,j_n}
 e^{\kappa\tau(\{i,j_1,\ldots,j_n\})}
             \norm{V_{n;i;j_1,\ldots,j_n}}.
\end{equation}
The coefficient norm is the multilinear operator norm of
the cell Banach spaces. All input cell labels are summed.

\begin{lemma}[Finite-range analytic vertices]\label{cv-local-lemma}
Suppose each output reads only labels within distance $R$,
at most $D_R$ labels, and $F$ is analytic and uniformly
bounded by $M_F$ on a global ball of radius $r_F$. If
$\norm{\mathscr R}\le C_P$, then
\begin{equation}\label{cv-local-majorant}
 v_n(\kappa)\le M'_F
       \left(\frac{eC_PD_Re^{\kappa R}}{r_F}\right)^n,
\end{equation}
with a fixed output-localization factor in $M'_F$.
If $F$ starts at degree two, the generating majorant is
analytic and quadratic at zero.
\end{lemma}

\begin{proof}
For normalized directions, multivariable Cauchy with radii
$r_F/(nC_P)$ bounds the derivative by
$M'_F(nC_P/r_F)^n$. Division by $n!$ and
$n^n/n!\le e^n$ gives the coefficient bound. There are at
most $D_R^n$ nonzero ordered tuples. Joining the root to
each input by a shortest path costs at most $nR$, proving
\eqref{cv-local-majorant}. This argument never counts
individual fine coordinates. On a smaller analytic margin
it also works at a nonzero fixed background, though the
constant and linear terms then need not vanish.
\end{proof}

Rooting at an input and summing the output and other inputs
has the same bound up to fixed factors, since the possible
outputs lie in a fixed neighborhood of that input.
For support-family substitution the exact geometric bound is
\begin{equation}\label{cv-tree-union}
 \tau\left(\{i\}\cup\bigcup_rX_r\right)
 \le\tau(\{i,j_1,\ldots,j_n\})
                    +\sum_r\tau(\{j_r\}\cup X_r).
\end{equation}
Join the graphs at their roots. The $j_r$ remain allowed as
unmarked auxiliary vertices when they are no longer terminals.
This constant-one inequality can fail for metric
minimum-spanning-tree length. If that convention is denoted
by $\ell_{\rm MST}$, a doubled traversal and shortcut give
\begin{equation}\label{cv-mst-conversion}
 \ell_{\rm MST}(S)\le2\tau(S).
\end{equation}
Thus a final weight $e^{\mu\ell_{\rm MST}}$ is covered by
a Steiner exponent at least $2\mu$. That spatial reserve
must be retained at the polymer interface.

\paragraph{Arbitrary fixed cardinality reserve.}
The concrete coefficients here are expanded at the fixed
flat background. They introduce no extra variable support.
For any prescribed $a\ge0$,
\begin{equation}\label{cv-cardinality}
 e^{a|\cup_rX_r|}\le\prod_re^{a|X_r|}.
\end{equation}
Consequently adding $e^{a|X|}$ to the rooted support norm
does not change the nonlinear vertex majorants. Linear
operators change the root, not the variable support. Only
the initial source $B_j$, with support $\{j\}$, costs
$e^a$. The source radius may shrink with $a$ while remaining
uniform in refinement and volume. This can pay a fixed
support-cardinality factor in a later polymer norm; small
spatial decay alone would not pay that factor. Hiding a
background-dependent minimizer in the Taylor coefficients
would require a new support argument and is not done here.

\subsection{The actual nonlinear vertices}
\label{cv-concrete-section}

\paragraph{The Wilson jet coefficient.}
For a positive plaquette use
$z=(a,b,c,e)=(A_\mu,A_\nu,\nabla_\mu A_\nu,
\nabla_\nu A_\mu)$ at its base point. The exact expressions are
\begin{equation}\label{cv-wilson-jet}
 \begin{split}
 P_\eta(z)&=e^{i\eta a}e^{i\eta b+i\eta^2c}
                 e^{-i\eta a-i\eta^2e}e^{-i\eta b},\\
 \mathcal L_\eta(z)&=-\frac12\tau\left[
       \left(\frac{P_\eta(z)-I}{\eta^2}\right)^2
                                      P_\eta(z)^{-1}\right].
 \end{split}
\end{equation}
The singularity at $\eta=0$ is removed before differentiating
vertices: $P_0=I$ and $\partial_\eta P_\eta|_0=0$.
Joint analyticity on a fixed jet ball gives uniform
coefficients. The bounded jet map and the $C^\alpha$
product estimate make
\[
 F_\eta^{\rm nl}(A)=f_\eta(SA):X_\eta\longrightarrow Y_\eta
\]
analytic on a uniform ball and quadratic at zero. Each
plaquette output reads a bounded physical neighborhood.
Lemma~\ref{cv-local-lemma} therefore proves its full
mixed-insertion bound uniformly in $\eta$.
The transpose $S^T$ stays outside this vertex. A uniform
bounded-field estimate for the unpropagated fine force
$S^TF_\eta^{\rm nl}$ is neither used nor claimed.

\paragraph{The Wilson energy density.}
Set
\begin{equation}\label{cv-cell-energy}
 \mathcal E_{i,\eta}(A)=\eta^4\sum_{p>0}
       \chi_i(x_p)\mathcal L_\eta((SA)_p).
\end{equation}
Then $\mathcal S_\eta(A)=\sum_i\mathcal E_{i,\eta}(A)$
exactly. Although a current cell contains $O(\eta^{-4})$
fine plaquettes,
\begin{equation}\label{cv-cell-mass}
 \eta^4\sum_{p>0}|\chi_i(x_p)|\le C_d.
\end{equation}
The fine differences are already bounded in the jet norm.
Thus every cell energy is a uniformly bounded analytic
scalar map on a fixed $\mathcal N_\alpha$ ball, with fixed
current dependency range. Its coefficients satisfy
\eqref{cv-local-majorant}. Replacing $\mathcal L_\eta$ by
its remainder after the quadratic Maxwell density gives
the same conclusion starting at degree three. Dropping
$\eta^4$, or treating all fine differences as unrelated
bounded coordinates, would lose this estimate.
The ordinary constrained scalar action is the same sum at
the corrected minimizer by the established gauge transfer.

\paragraph{The full corrected constraint.}
At input spacing $h$ and output spacing $Lh$,
\begin{equation}\label{cv-one-step}
 \begin{split}
 \mathcal F_h(A)&=(Lh)^{-1}F_L(hA)=Q_hA+\mathcal R_h(A),\\
 \norm{\mathcal R_h(A)}_\infty&\le Ch\norm A_\infty^2,
 \qquad
 \norm{D\mathcal R_h(A)}_{\infty\to\infty}
       \le Dh\norm A_\infty.
 \end{split}
\end{equation}
For $h_t=L^{t-j}$, $\sum_{t<j}h_t\le(L-1)^{-1}$.
The proved Riccati bound keeps every intermediate physical
field in a common ball. Hence the composed remainder
$\mathcal R$ is uniformly analytic, quadratic, and has
physical dependency radius at most
$r_L\sum_th_t\le R_*$.
It follows that $\mathcal R:X_\eta\to E$ satisfies
Lemma~\ref{cv-local-lemma}. The scale sum is already included
in this uniform bound; higher coefficients are obtained by
Cauchy in that ball, not by an uncontrolled count of
multilevel diagrams.

\paragraph{The physical-adjoint constraint vertex.}
Define
\begin{equation}\label{cv-adjoint-vertex}
 \mathcal V_\eta(A,m)=D\mathcal R(A)^Tm:
                X_\eta\times E\longrightarrow Z_\eta.
\end{equation}
Ordinary holomorphy of $\mathcal R$ alone would not control
this physical transpose. The proved row estimate does:
at one step,
\[
 \norm{D\mathcal R_h(A)^T}_{\infty\to\infty}
       \le D_*h\norm A_\infty,\qquad
 D_*=N_LK_0L^{d-1}.
\]
The factor $L^d$ in the physical transpose is included here.
Both $Q_h$ and $Q_h^T$ contract the bounded norm.
Transposing the composition and using the intermediate
field bound proves
\begin{equation}\label{cv-adjoint-bound}
 \norm{D\mathcal R(A)^T}_{\infty\to\infty}
    \le \exp\left(\frac{2D_*\norm A_\infty}{L-1}\right)-1.
\end{equation}
The map is holomorphic into this physical absolute-row
operator space. Thus \eqref{cv-adjoint-vertex} is jointly
analytic and $O(\norm A_\infty\norm m_\infty)$.
Its output at a fine bond reads coarse $m$ within $R_*$,
and background $A$ within $2R_*$, because each relevant
coarse equation contains that bond and reads its own
neighborhood. Lemma~\ref{cv-local-lemma} applies with
total degree starting at two. Cauchy is taken in the
physical-row operator norm; it is not applied entrywise
and followed by a count of fine coefficients.
The output is only required to be bounded: it enters
$\mathcal C$ or $H^T$, and no derivative of $m$ is taken.

\subsection{Summable linear blocks, including the derivative kernels}
\label{cv-linear-section}

The established complete-field estimates, for every positive
weight with $w(i)/w(j)\le e^{\kappa_0d(i,j)}$, are
\begin{equation}\label{cv-linear-inputs}
 \begin{gathered}
 H:E_w\to X_{\eta,w},\qquad
 \mathcal C:Z_{\eta,w}\to X_{\eta,w},\qquad
 \mathcal CS^T:Y_{\eta,w}\to X_{\eta,w},\\
 H^T:Z_{\eta,w}\to E_w,\qquad
 (SH)^T:Z_{\eta,w}\to E_w,\qquad q:E_w\to E_w.
 \end{gathered}
\end{equation}
Restricting the fifth input to $Y_\eta$ is sufficient.
The third estimate has the same $\alpha$ at input and
output; a positive H\"older loss would not suffice.
These are the constrained Schauder and extension bounds
already used in the weighted response proof. They do not
require the gauge projection $\Pi$ to act on arbitrary
bounded fields.

For any such $T$, form $T_{ij}=P_iTI_j$. Given a supported
input $u_j$, use $w_j(n)=e^{\kappa_0d(n,j)}$. Its weighted
norm is bounded by its unweighted norm times a fixed collar
factor. The complete weighted estimate gives
\begin{equation}\label{cv-block-envelope}
 \norm{T_{ij}}\le C_Te^{-\kappa_0d(i,j)}.
\end{equation}
This is an operator norm between cell Banach spaces, not an
absolute integral estimate for a singular differentiated
fine-coordinate kernel. In particular its diagonal block is
controlled by the equal-exponent Schauder bound.
For every $0<\kappa<\kappa_0$, lattice growth then gives
\begin{equation}\label{cv-block-schur}
 \max\left\{
  \sup_i\sum_je^{\kappa d(i,j)}\norm{T_{ij}},\,
  \sup_j\sum_ie^{\kappa d(i,j)}\norm{T_{ij}}\right\}
                \le C_{T,\kappa}.
\end{equation}
Both sums follow from the common envelope, without any
self-adjointness assumption in these Banach norms.
Periodic-distance sums have the same uniform upper bound.

The decomposition is exact; for example,
\begin{equation}\label{cv-exact-linear-splitting}
 \mathcal CS^Tf=\sum_j\mathcal CS^T(\chi_jf).
\end{equation}
Every $\chi_jf$ is a complete $C^\alpha$ field. Smooth
cutoff derivatives are included in $S^T(\chi_jf)$ and
controlled by the same estimate; they cancel in the sum.
No sharp splitting is used. The maps $H^T$ and $(SH)^T$
likewise retain their physical transpose normalization.
Output first differences and their H\"older seminorms are
already in $X_\eta$, while source derivatives occur only
through $\mathcal CS^T$ and $(SH)^T$. Separate absolute
fine-coordinate derivative bounds are unnecessary for
\eqref{cv-vertex-norm}.

\subsection{An exact simultaneous fixed point using these vertices}
\label{cv-fixed-point-section}

The corrected constraint is $\mathcal F(A)=QA+\mathcal R(A)$.
Let $m$ be the multiplier of Notebook15 and write
$\mathcal V_\eta(A,m)=D\mathcal R(A)^Tm$.
The exact stationarity equations are equivalent near zero to
\begin{equation}\label{cv-fixed-point}
 \begin{split}
 A={}&HB-H\mathcal R(A)
          -\mathcal CS^Tf_\eta(SA)+\mathcal C\mathcal V_\eta(A,m),\\
 m={}&qB-q\mathcal R(A)
          +(SH)^Tf_\eta(SA)-H^T\mathcal V_\eta(A,m).
 \end{split}
\end{equation}
Indeed the second equation rearranges to
\[
 (I+H^TD\mathcal R(A)^T)m
        =q[B-\mathcal R(A)]+H^TN(A),
\]
and the first is $A=H[B-\mathcal R(A)]+\xi$ with
$\xi=-\mathcal CN(A)+\mathcal CD\mathcal R(A)^Tm$.
These are the proved constraint chart and reduced
stationarity identities. Conversely the exact constraints
of $\mathcal C$ put $\xi$ in the fixed flat slice, and the
small multiplier inverse recovers the earlier formulation.
No extra small critical points are introduced.

All nonlinear primitives in \eqref{cv-fixed-point} start at
degree two: $\mathcal R(A)$, $f_\eta(SA)$, and
$D\mathcal R(A)^Tm$. Their mixed-insertion majorants are
analytic and uniform by the preceding proof, while their
linear propagators satisfy \eqref{cv-block-schur}.
Apply $P$ to the equations and reconstruct with $\mathscr R$.
A fixed point in compatible cell components reconstructs
exactly to \eqref{cv-fixed-point}; the overlap constants are
absorbed by reducing the small radius. This supplies the
concrete inputs for the separate rooted support-family
fixed-point theorem, including any fixed cardinality reserve.
Tree summation, canonical grouping of supports and the final
scalar-action expansion remain conclusions of that theorem,
not substitutes for its proof.

\subsection{An obstruction to using the full ordinary adjoint}
\label{cv-ordinary-section}

The full ordinary rooted average is analytic from small
complete physical connections to bounded current data, with
fixed physical dependency range, by the ordinary trajectory
theorem. It therefore has $X_\eta\to E$ vertex bounds.
Its physical adjoint on arbitrary bounded current fields is
different and is not uniformly bounded through refinement.

Let $Q_{\rm ord}$ be its flat derivative in physical
potentials. Exact ordinary gauge covariance gives
\begin{equation}\label{cv-ordinary-gradient}
 Q_{\rm ord}D_\eta\lambda
          =D_1(\lambda|_{\text{final centers}}),
\end{equation}
with the same sign on both sides for either consistent
gauge convention. On a sufficiently large torus, take a
unit color vector $\lambda$ at one final-center fine site
and zero at every other fine site. Its fine gradient is
nonzero on $2d$ bonds, each of magnitude $\eta^{-1}$, so
\begin{equation}\label{cv-point-gradient-norms}
 \norm{D_\eta\lambda}_{2,\eta}
       =\sqrt{2d}\,\eta^{d/2-1},\qquad
 \norm{D_\eta\lambda}_{1,\eta}=2d\eta^{d-1}.
\end{equation}
The current gradient in \eqref{cv-ordinary-gradient} has
magnitude one on $2d$ bonds. Therefore
\begin{equation}\label{cv-ordinary-hilbert-obstruction}
 \norm{Q_{\rm ord}}_{L^2_\eta\to L^2_1}
                     \ge\eta^{1-d/2}.
\end{equation}
Choose $m$ to equal that current gradient, so
$\norm m_\infty=1$ and
$\langle Q_{\rm ord}D_\eta\lambda,m\rangle_1=2d$.
Physical transposition and
\eqref{cv-point-gradient-norms} imply
\begin{equation}\label{cv-ordinary-adjoint-obstruction}
 \norm{Q_{\rm ord}^T}_{\infty\to\infty}\ge\eta^{1-d}.
\end{equation}
In dimension four these are $\eta^{-1}$ and $\eta^{-3}$.
The complete $\mathcal N_\alpha$ norm of the test gradient
is large, so there is no contradiction with
$Q_{\rm ord}:X_\eta\to E$ boundedness.

The corrected flat derivative instead satisfies
$Q_M^{\rm st}D_\eta=D_1Q_0$ with the scalar block mean,
and its physical transpose contracts the bounded norm.
This is why \eqref{cv-fixed-point} uses the corrected map.
Possible cancellations in a specially projected ordinary
adjoint are not assumed. The ordinary current-step pivot
chart of Notebook19 is unaffected: that ratio is the fixed
$L$, so its constants may depend on $L$. The obstruction
concerns unbounded total ratio $M$.

\subsection{Scope of the concrete input}
\label{cv-scope}

The proved inputs are uniformly analytic finite-current-range
Wilson jet and energy vertices, corrected-constraint and
physical-adjoint constraint vertices, summed mixed-insertion
norms, and exponentially summable linear cell blocks including
the derivatives used in the simultaneous equations. The
one-weight theorem is used only to form linear block
envelopes; the nonlinear coefficient proof really sums all
input labels.

No raw fine-coordinate absolute Taylor norm without physical
normalization or regularity is asserted. No ordinary rooted
adjoint bound, independent-current source extension,
mixed-history theorem or quantum integration result is
inferred. Backgrounds remain in the homogeneous small global
potential branch. The scalar gauge transfer, rather than a
false identification of ordinary and corrected derivatives,
matches its classical action to the ordinary branch.

\section{Strict support localization of the homogeneous classical action}
\label{cl-section}

Fix odd $L$, dimension four, $M=L^j$, $\eta=M^{-1}$ and
$0<\alpha<1$. Work first on compatible finite periodic regulators.
Constants below are independent of $M$ and volume. They may depend on
$L$, color, $\alpha$, the requested positive spatial exponent and an
arbitrarily prescribed finite support-cardinality reserve. The product
background radius may shrink with these fixed choices. Fine transposes
retain the physical measure $\eta^4$.

The concrete cell-vertex estimates preceding this section supply the
exact simultaneous equations
\begin{align}
 A&=HB-HR(A)-\mathcal C S^T f_\eta(SA)
                         +\mathcal C V(A,m),\notag\\
 m&=qB-qR(A)+(SH)^T f_\eta(SA)-H^T V(A,m),
 &V(A,m)&=DR(A)^Tm.                                      \label{cl-fixed-point}
\end{align}
Here $\mathcal F(A)=QA+R(A)$ is the full corrected fixed-reference
constraint, $H$ is the matched flat minimizing extension,
$\mathcal C$ the flat constrained inverse, $q=q_M$, and $S$ the local
first-difference jet. All three nonlinear primitives start at total
degree two. Rearranging the second equation gives the multiplier inverse
of the constrained stationary equation; the first then gives its exact
constraint chart. Thus \eqref{cl-fixed-point} is equivalent to the
previously proved small classical branch. It uses the physical adjoint
of the corrected constraint, not the nonuniform adjoint of the full
ordinary rooted average. The scalar action equality supplied by the
ordinary-constraint gauge transfer will be used after solving these
equations.

\subsection{Geometry and the support-family norm}
\label{cl-geometry-section}

Let $d$ be current-cell graph distance. For a nonempty finite set $X$,
let $\tau(X)$ be the minimum number of edges in a connected graph
containing $X$, allowing unmarked intermediate vertices. This graph
Steiner length differs from the metric minimum spanning-tree length
$\ell$ used in the Gaussian expansion, but
\begin{equation}
 \tau(X)\le\ell(X)\le2\tau(X),\qquad |X|\le\tau(X)+1.       \label{cl-metric-comparison}
\end{equation}
The second upper bound follows by traversing a connected minimizing
tree and shortcutting its tour in the graph metric. The crucial rooted
inequalities are
\begin{align}
 \tau(\{i\}\cup X)&\le d(i,j)+\tau(\{j\}\cup X),\notag\\
 \tau\left(\{i\}\cup\bigcup_{r=1}^nX_r\right)
 &\le\tau(\{i,j_1,\ldots,j_n\})
                         +\sum_{r=1}^n\tau(\{j_r\}\cup X_r).
                                                               \label{cl-root-gluing}
\end{align}
The auxiliary roots $j_r$ need not be final support labels. This is why
the support algebra uses $\tau$: the analogous inequality for metric
spanning trees with deleted auxiliary vertices need not hold. All these
statements remain valid on periodic graphs.

Use the smooth fixed-physical-scale partition of unity and the complete
cell spaces from the concrete vertex construction. Their field types
are $\mathcal N_\alpha$, $C^\alpha$, bounded fine fields, and
fixed-dimensional current fields. Localization $P_i$ and reconstruction
$\mathcal R$ are uniformly bounded, and $\mathcal R P=I$.
In particular, input pieces to $\mathcal C S^T$ remain complete
$C^\alpha$ fields; no sharp cut is differentiated. For each relevant
linear propagator $T$, the proved weighted complete-field estimate gives
\begin{equation}
 \|T_{ij}\|\le C_Te^{-\kappa_0d(i,j)},\qquad
 L_T(\kappa):=\sup_i\sum_j e^{\kappa d(i,j)}\|T_{ij}\|<\infty
 \quad(0<\kappa<\kappa_0).                                  \label{cl-linear-blocks}
\end{equation}
The diagonal block of $\mathcal C S^T$ is controlled by the proved
equal-exponent Schauder estimate, rather than an absolute singular-kernel
integral. These block estimates are uniform in refinement and volume.

On the original current-background product ball of radius $b$, a support
family is a collection $u_{i;X}(B|_X)$ of bounded holomorphic functions
with values in the cell space at $i$. The support $X$ is nonempty and
records original background cells. For $a\ge0$ put
\begin{equation}
 \|u\|_{\kappa,a;b}
 :=\sup_i\sum_X e^{\kappa\tau(\{i\}\cup X)+a|X|}
       \sup_{\|B|_X\|_\infty<b}\|u_{i;X}(B|_X)\|.             \label{cl-family-norm}
\end{equation}
For several field types, take their maximum in each coefficient norm
\emph{before} the support sum: the joint family norm is
$\sup_i\sum_X e^{\kappa\tau(\{i\}\cup X)+a|X|}
\max_t\|u^t_{i;X}\|$, with the displayed background supremum understood.
It is not the maximum of the separate type-family norms. Taylor
operator norms below use this same joint cell space with its maximum
norm. A bound for a joint linear source may be taken as the sum of its
component block bounds. These are fixed finite sums.
This is a Banach space. Terms with the same support are added in that
slot; evaluation $\sum_Xu_{i;X}(B|_X)$ converges absolutely. All
coefficients remain functions of the original variables. No nonlocal
background-dependent minimizer is hidden in their definition.

\begin{lemma}[Support substitution]\label{cl-substitution}
A linear propagator obeying \eqref{cl-linear-blocks} has norm at most
$L_T(\kappa)$ in \eqref{cl-family-norm}. Suppose a nonlinear cell map has
Taylor coefficients, including $1/n!$, whose mixed-insertion norms obey
\begin{align}
 v_n&:=\sup_i\sum_{j_1,\ldots,j_n}
 e^{\kappa\tau(\{i,j_1,\ldots,j_n\})}
       \|N_{n;i;j_1,\ldots,j_n}\|,\notag\\
 \Phi(t)&:=\sum_{n\ge2}v_nt^n<\infty\quad(0\le t<t_0).
                                                               \label{cl-vertex-majorant}
\end{align}
Assign a substituted term the union of its input supports. Then
\begin{align}
 \|N(u)\|_{\kappa,a;b}&\le\Phi(\|u\|_{\kappa,a;b}),\notag\\
 \|N(u)-N(v)\|_{\kappa,a;b}&\le\Phi'(r)\|u-v\|_{\kappa,a;b}
 \quad(\|u\|,\|v\|\le r<t_0).                              \label{cl-substitution-bound}
\end{align}
The majorant is independent of the prescribed cardinality reserve $a$.
\end{lemma}

\begin{proof}
For the linear assertion attach a path from the input root $j$ to the
output root $i$ and use the first inequality in
\eqref{cl-root-gluing}. For a nonlinear term use the second inequality
and
$|\bigcup_rX_r|\le\sum_r|X_r|$.
The absolute sum over all supports and ordered input-cell slots then
factorizes into the product of their family norms, leaving $v_n$.
Summing over $n$ proves the first bound. Telescoping an $n$-linear
expression gives the factor $n r^{n-1}$ and proves the second.
Repeated supports, repeated labels and overlapping supports are all
included in these sums. A Taylor coefficient at the fixed flat
background introduces no additional variable support, so it adds no
cardinality cost.
\end{proof}

The concrete finite-range analytic vertices have
\begin{equation}
 v_n\le M_1
       \left(\frac{eC_PD_Re^{\kappa R}}{r_0}\right)^n.        \label{cl-local-cauchy}
\end{equation}
Here $R$ is a uniform current-cell dependency radius, $D_R$ the fixed
number of admissible input cells in one slot, $C_P$ the reconstruction
norm, and $r_0$ a uniform analytic radius. Indeed multivariable Cauchy
with radius $r_0/(nC_P)$ bounds the derivative by
$M_1(nC_P/r_0)^n$; division by $n!$ and $n^n/n!\le e^n$ bounds its
Taylor coefficient. There are at most $D_R^n$ nonzero ordered cell
tuples, with a star of length at most $nR$ through the output.
This counts current cells, not $M^4$ independent fine coordinates.
For $V(A,m)$ it uses the previously proved holomorphy in the physical
absolute-row norm of $DR(A)^T$; ordinary bounded-field holomorphy of
$R$ alone would not imply that input.

\subsection{A convergent support-valued stationary solution}
\label{cl-contraction-section}

\begin{theorem}[Uniform strict-support fixed point]
\label{cl-fixed-point-theorem}
For every fixed finite $a\ge0$ and sufficiently small positive
$\kappa$, there is $b_*>0$, independent of $M$ and volume, such that
the solution of \eqref{cl-fixed-point} has support families satisfying
\begin{align}
 \|(A,m)\|_{\kappa,a;b}&\le C_a b,\notag\\
 \|(A-HB,m-qB)\|_{\kappa,a;b}&\le C_a b^2
 \qquad(0<b\le b_*).                                      \label{cl-solution-bound}
\end{align}
Every coefficient of the second line has Taylor degree at least two.
\end{theorem}

\begin{proof}
Apply the preceding lemma to the three nonlinear primitives in
\eqref{cl-fixed-point} and then their linear propagators. Their combined
positive coefficient majorant can be bounded coefficientwise by
\begin{equation}
 \Phi(t)\le\frac{v(ct)^2}{1-ct},\qquad
 \Phi'(t)\le8vc^2t\quad(ct\le1/2),                          \label{cl-majorant-constants}
\end{equation}
after increasing fixed $v,c$. The singleton source
$B_{i;\{i\}}=B_i$ has norm $e^ab$. If $L_0$ bounds the linear map
$B\mapsto(HB,qB)$, choose
\begin{equation}
 r_*\le\min\{(2c)^{-1},(16vc^2)^{-1}\},\qquad
 b_*\le\frac{r_*}{2L_0e^a}.                                \label{cl-radius}
\end{equation}
Also shrink $r_*$ so that evaluation and reconstruction stay in the
previously established classical branch's uniqueness ball; their norms
are uniformly bounded. The support map is a contraction with constant
at most $1/2$ on the ball of radius $2L_0e^ab$. This proves the first
estimate. Applying the quadratic majorant to its solution proves the
second. Degree induction in Picard iteration proves the vanishing
orders coefficientwise.

Evaluation is continuous and respects every absolutely convergent
substitution. The evaluated solution therefore solves the exact
equations \eqref{cl-fixed-point}. The compatible cell image is preserved
by their right-hand side, and uniqueness in the common small domain
identifies the result with the proved classical minimizer and
multiplier. Picard iteration also defines an absolutely convergent
rooted-tree expansion: local Taylor coefficients are vertices, flat
cell propagators are edges, and the leaves carry the original $B_i$.
All depths, orders and repeated supports are already summed by the
Banach contraction. No additional factor counting the number of RG
levels occurs.
\end{proof}

\subsection{Scalar action and symmetric support Hessians}
\label{cl-action-section}

The established envelope identity is
\begin{equation}
 DS_M(B)=m(B),\qquad S_M(0)=0,\qquad D^2S_M(0)=q_M.          \label{cl-envelope-input}
\end{equation}
The exact gauge transfer identifies this scalar with the ordinary
homogeneous constrained action. Write $n=m-qB$. On a finite regulator,
\begin{equation}
 S_M(B)-\tfrac12\langle B,qB\rangle
    =\int_0^1\sum_i\langle B_i,n_i(tB)\rangle\,dt.           \label{cl-envelope}
\end{equation}
The current pairing has unit cell volume and fixed fiber dimension.
Define strictly supported scalar coefficients by
\begin{equation}
 W_Y(B|_Y)=
 \sum_{\{i\}\cup X=Y}\int_0^1
       \langle B_i,n_{i;X}(tB|_X)\rangle\,dt.                \label{cl-action-coefficients}
\end{equation}
They are holomorphic and begin at order three.

\begin{proposition}[Strict scalar action localization]
\label{cl-action-theorem}
For $0<\mu<\kappa$ and every prescribed finite $a\ge0$, on a fixed
small product ball,
\begin{align}
 S_M(B)&=\tfrac12\langle B,q_MB\rangle+\sum_{Y\ne\varnothing}W_Y(B|_Y),
                                                               \label{cl-action-expansion}\\
 \sup_b\sum_{Y\ni b}e^{\mu\tau(Y)+a|Y|}\|W_Y\|_r
       &\le C_{\mu,a}r^3.                                    \label{cl-action-norm}
\end{align}
Here $\|W_Y\|_r$ is the supremum on the complex product ball of radius
$r$. Each Hessian $B_Y=D^2W_Y$, extended by zero outside cell rows and
columns in $Y$, is individually complex symmetric and depends only on
$B|_Y$. With a fixed cardinality spare $\epsilon>0$,
\begin{equation}
 \sup_i\sum_{Y\ni i}\sum_{j\in Y}
 e^{\mu\tau(Y)+(a-\epsilon)|Y|}
            \|(B_Y)_{ij}\|_r\le C r.                       \label{cl-hessian-support}
\end{equation}
\end{proposition}

\begin{proof}
For $b\in Y=\{i\}\cup X$, one has $d(i,b)\le\tau(Y)$ and
$|Y|\le|X|+1$. Thus
\begin{equation}
 e^{\mu\tau(Y)+a|Y|}
 \le e^a e^{\kappa\tau(Y)+a|X|}
                   e^{-(\kappa-\mu)d(i,b)}.                \label{cl-reroot}
\end{equation}
Summing first over $X$ and then $i$, and using
$\sup_b\sum_i e^{-\delta d(i,b)}\le\sum_{x\in\mathbb Z^4}e^{-\delta|x|}$,
proves \eqref{cl-action-norm} from \eqref{cl-solution-bound} and
\eqref{cl-action-coefficients}. This is an anchored sum over all supports
containing an arbitrary cell, not only the original output root.

Use the action estimate on radius $2r$ and take derivatives on radius
$r$. Cauchy and polarization give each cell block the dimension-free
bound
\begin{equation}
 \|(D_iD_jW_Y)\|_r\le\frac{2e^2}{r^2}\|W_Y\|_{2r}.         \label{cl-hessian-cauchy}
\end{equation}
The row sum contributes at most $|Y|$, absorbed by
$\sup_{n\ge1}ne^{-\epsilon n}<\infty$ in the reserve. This proves
\eqref{cl-hessian-support} and justifies termwise differentiation.
Scalar holomorphy gives symmetry of each coefficient. Rows, columns and
background support are all in the same $Y$ by construction.
\end{proof}

Taking $\mu\ge2(\kappa_{\rm G}+2m_{\rm G})$ converts
\eqref{cl-hessian-support} into the strict precision-support norm with
metric-tree weight $e^{(\kappa_{\rm G}+2m_{\rm G})\ell(Y)}$ used by the
Gaussian density algebra. The Gaussian exponents can be chosen positive
and sufficiently small compared with the available flat decay. The
smallness of the precision norm is then obtained by shrinking the
background radius, not by an unsupported claim about a large original
source domain.

\subsection{The actual current pivot chart and retained precision}
\label{cl-chart-section}

Use the fixed-$L$ ordinary pivot and coordinate maps already identified
for the physical current step. With $a=\log V/i$, set
\begin{equation}
 F_{\rm cur}(a,z)=\frac1i\log\!\left(e^{i\Phi_V(z)}e^{ia}\right),
 \qquad F_{\rm cur}(a,0)=a.                                  \label{cl-current-chart}
\end{equation}
They are holomorphic on a fixed joint product neighborhood, have a fixed
current-cell dependency radius $R_c$, and satisfy
\begin{equation}
 \|F_{\rm cur}(a,z)\|_\infty
             \le C_L(\|a\|_\infty+\|z\|_\infty).            \label{cl-current-chart-bound}
\end{equation}
Let $G_{\rm cur}(a,z)$ be the actual local exponential gauge action on
$e^{i\Phi_V(z)}V$, not a replacement relative quadratic penalty. Its
full analytic expression is finite range at fixed $L$. At the physical
rooted-contour background its quadratic form is the matched gauge term;
away from that locus its actual local curvature terms remain included.
Put
\begin{equation}
 T_M(a,z)=S_M(F_{\rm cur}(a,z))+G_{\rm cur}(a,z).             \label{cl-pulled-action}
\end{equation}

The scalar family for $S_M$ includes its quadratic baseline: split
$\frac12\langle B,qB\rangle$ into pair-supported scalar coefficients.
Exponential locality of $q$ gives their anchored support norm of order
$r^2$ with any fixed cardinality reserve and a sufficiently small
positive spatial exponent. Adding \eqref{cl-action-coefficients}
therefore gives a scalar support family for the entire $S_M$.

For a support $Y$ let $\Gamma Y$ be the union of the radius-$R_c$
dependency neighborhoods of its cells. Fixed constants $N_c,c_c$ obey
\begin{equation}
 |\Gamma Y|\le N_c|Y|,\qquad
 \tau(\Gamma Y)\le\tau(Y)+c_c|Y|.                           \label{cl-support-enlargement}
\end{equation}
Attach a fixed finite neighborhood tree at each point of a minimizing
connected graph for $Y$ to prove the second inequality. Periodic
identifications only reduce this count. Hence
\begin{equation}
 e^{\sigma\tau(\Gamma Y)+a_{\rm out}|\Gamma Y|}
 \le e^{\sigma\tau(Y)+(\sigma c_c+a_{\rm out}N_c)|Y|}.       \label{cl-composition-weight}
\end{equation}
Choose the input cardinality reserve to pay this fixed cost, and shrink
the joint radius so that \eqref{cl-current-chart-bound} remains in the
source-action ball. Substitute the chart in each scalar coefficient and
assign support $X=\Gamma Y$. A root in $\Gamma Y$ can be moved to a
point of $Y$ in at most a fixed number of ways. Adding the local gauge
coefficients gives the exact expansion
\begin{equation}
 T_M(a,z)=\sum_X T_X(a|_X,z|_X),                             \label{cl-pulled-support-expansion}
\end{equation}
with a uniform anchored scalar norm on the fixed joint product domain
and any prescribed fixed output cardinality reserve. Both the original
background and retained-coordinate labels remain explicit in $X$.

\begin{corollary}[Strict support of the physical retained precision]
\label{cl-precision-theorem}
Define
\begin{equation}
 A_X(a|_X)=D_z^2T_X(a|_X,0)-D_z^2T_X(0,0).                   \label{cl-precision-coefficients}
\end{equation}
These are individually complex-symmetric retained matrices with rows,
columns and background dependence in $X$. For some $\sigma>0$ and all
sufficiently small $b$,
\begin{align}
 A_M(V)&=A_M(1)+\sum_X A_X((\log V/i)|_X),                  \label{cl-precision-expansion}\\
 \sup_i\sum_{X\ni i}\sum_{j\in X}
 e^{\sigma\tau(X)}\|(A_X)_{ij}\|_b&\le C_L b.             \label{cl-precision-norm}
\end{align}
The baseline is exactly the matched flat retained precision
$A_M(1)=K_L^T(q_M+T_L^TT_L)K_L$.
\end{corollary}

\begin{proof}
Apply Cauchy with fixed margins to the joint domain of
\eqref{cl-pulled-support-expansion}, twice in $z$ and once in $a$.
Integrating the last derivative along $ta$ gives $O(b)$ in each
coefficient. Its row sum costs $|X|$, paid by an available cardinality
reserve as in \eqref{cl-hessian-support}. The baseline sum is the flat
Hessian of the exact pulled action and hence the identified retained
precision. The matrices are Hessians of scalar coefficients, which
proves their individual symmetry and supports.
\end{proof}

This argument retains both nonlinear constraint curvature and
exponential-coordinate curvature. It does not identify a freely chosen
independent source current with the physical Hessian. The full product
background extension is an analytic Hessian family; only the physical
background is asserted to be stationary. Choosing
$2(\kappa_{\rm G}+2m_{\rm G})<\sigma$ and then shrinking $b$ makes the
strict norm as small as the stated Gaussian support algebra requires.
The baseline covariance is background independent and has the uniform
spectrum and decay proved in the flat theory.

\subsection{Localized cubic interaction and Gaussian-cell regrouping}
\label{cl-interaction-section}

For $\|z\|_\infty\le R$ and $gR$ inside a fixed analytic margin define
\begin{align}
 U_X(a;z)=-g^{-2}\bigl[&T_X(a,gz)-T_X(a,0)
          -gD_zT_X(a,0)[z]\notag\\
          &-\tfrac12g^2D_z^2T_X(a,0)[z,z]\bigr].             \label{cl-cubic-coefficients}
\end{align}
The sign corresponds to $\exp(-g^{-2}T_M)$.
Taylor's integral remainder, or one-variable Cauchy along $z/R$, yields
\begin{equation}
 \sup_i\sum_{X\ni i}e^{\sigma\tau(X)+a_{\rm out}|X|}
 \sup_{\|a\|<b,\,\|z\|\le R}|U_X(a;z)|\le CgR^3.          \label{cl-cubic-norm}
\end{equation}
No factor $|X|^3$ is needed: the full direction $z|_X$ has product norm
at most $R$. At a physical stationary background the sum of the linear
terms in \eqref{cl-cubic-coefficients} vanishes. This cancellation is
not asserted for every point of the analytic product-background domain.

For the fixed retained product cutoff set
$\chi_X=\prod_{i\in X}\chi_i$ and
$w_X=\chi_X(e^{U_X}-1)$. Then
$v_X:=\sup|w_X|\le e^{\|U_X\|}\|U_X\|$.
Take $a_{\rm out}>4+\log\widehat K$ with a positive spare, where
$\widehat K$ is the chosen fixed Gaussian density factor, and choose
$2\kappa_{\rm G}\le\sigma$. Equations
\eqref{cl-metric-comparison} and \eqref{cl-cubic-norm} give
\begin{equation}
 \sup_i\sum_{X\ni i}v_X\widehat K^{|X|}
                   e^{4|X|+\kappa_{\rm G}\ell(X)}
        \le C e^{CgR^3}gR^3.                                \label{cl-polymer-input}
\end{equation}
Thus this is the required localized interaction norm, not only an
extensive unweighted cubic remainder. A logarithmically growing $R$ as
$g\downarrow0$ is compatible with both the analytic restriction on $gR$
and smallness of $gR^3$. The separate Gaussian bad-block tail and the
explicit density-algebra smallness conditions must still be imposed.
No all-bond nonlinear sharp cutoff is identified with the retained
product window by this argument.

The support algebra above uses unit current cells. The Gaussian theorem
may use one cell per fixed $L$-block, with
$D=d(L^d-1)\dim\mathfrak g$ retained real coordinates in the usual pivot
convention. Let $\pi$ map a unit current cell to its $L$-block and regroup
each support $Y$ to $\pi(Y)$. Then
\begin{equation}
 \tau_{\rm block}(\pi(Y))\le\tau_{\rm unit}(Y),\qquad
 |\pi(Y)|\le|Y|.                                           \label{cl-block-regrouping}
\end{equation}
Every unit edge either stays in one block or crosses to an adjacent
block. A block-anchored support sum is bounded by the sum of the unit
anchored sums over its at most $L^d$ cells. For matrices, bound a block
matrix norm by the sum of its constituent unit-cell blocks, then sum the
at most $L^d$ row cells. Thus scalar and absolute-row support norms
change by only fixed $L$-dependent factors. Maximum-cell and Euclidean
block norms likewise differ by fixed factors; reduce a background or
cutoff radius accordingly. No dimension depending on $M$ is introduced.
The Gaussian density constants in \eqref{cl-polymer-input} can therefore
use the fixed retained dimension $D$ after this explicit regrouping.

\subsection{Scope of the classical input}
\label{cl-scope}

The construction expands about a fixed flat baseline in the original
current variables. The arbitrary finite cardinality reserve is paid by
a smaller uniform product radius; spatial decay alone would not pay
the Gaussian factor $\widehat K^{|X|}e^{4|X|}$. Fine dimension is removed
by complete cell norms, the physical measure, and the actual
physical-adjoint row bound. The scalar ordinary action is identified by
the proved gauge transfer, while its current pivot chart has fixed
ratio $L$.

The conclusions supply strict-support precision and classical
interaction inputs for the homogeneous classical step at every depth
$M$, on the indicated small branch. They do not establish closure of
inherited quantum polymer terms, arbitrary independent-current or
mixed-history domains, a running-coupling theorem, or an interacting
continuum limit. Finite compatible periodic regulators are primary;
the graph estimates are uniform, but source chart incidence still uses
the admissible current tori assumed in its proof. No thermodynamic
consistency or separate gauge invariance of an arbitrarily regrouped
individual coefficient is asserted.

% No preamble. Follows cv- concrete inputs and cl- support localization.
% Unique prefix ci-.
\section{One integration of the classical action at every earlier depth}
\label{ci-section}

The support theorem upgrades the bare calculation of Notebook 20.
Its exact scope matters: $S_M$ is the classical action obtained
by minimizing the original Wilson action through the earlier
ratio $M$. The following integral uses $S_M$ and the source
gauge and measure factors. It has no previously generated
quantum correction $E_k$. Thus the result is uniform through
classical depth, not a claim that the quantum induction is
already closed.

Let $\Phi_V$ be Notebook 19's exact retained chart, let
$a=\log V/i$, and put
\begin{equation}\label{ci-action}
 \mathcal T_{M,V}(x)=
 S_M\bigl(\log(e^{i\Phi_V(x)}V)/i\bigr)
                    +\mathcal G(e^{i\Phi_V(x)}V).
\end{equation}
The full product-domain Hessian is
$A_M(V)=D_x^2\mathcal T_{M,V}(0)$, including the absolute
gauge function away from its zero slice. Its flat baseline is
$A_{M,0}=K_L^T(q_M+T_L^TT_L)K_L$.
The preceding support construction gives local scalar
coefficients $\mathcal T_{M,Y}$ and symmetric precision
coefficients with all their retained rows, columns and
background dependence in $Y$.

Group unit current cells into fixed $L$-blocks, assigning their
retained positive bonds to them. The Gaussian cell dimension
is then $D=d(L^d-1)\dim\mathfrak g$, independent of $M$ and
volume. Any bounded change of support convention and anchor
multiplicity in this grouping costs constants depending only
on $L$. After the spatial-exponent and cardinality reserves
proved above, the resulting coefficients have exactly the
weighted support norms of Notebook 18.

Define the logarithmic interaction $v_{M,Y}$ by subtracting
the constant, linear and quadratic terms from each
$-g^{-2}\mathcal T_{M,Y}(a,gz)$ and adding the normalized
straightening-Jacobian and Haar logarithms of Notebook 19.
The gauge Taylor remainder is already included in
$\mathcal T_{M,Y}$ and is not counted twice. There is a
common small product background ball such that, for every
fixed required support-cardinality weight and a sufficiently
small spatial exponent,
\begin{align}
 \nu_M(b)&\le C_L b,\label{ci-inputs}\\
 \sup_i\sum_{Y\ni i}e^{a_0|Y|+\kappa\ell(Y)}
                      \|v_{M,Y}\|_{b,R}
       &\le C_{L,a_0,\kappa}gR^3,
       \qquad R\ge1,\quad gR<r_L.\notag
\end{align}
Here $\nu_M$ is the precision support norm, not just the
norm of the summed matrix, and $\ell$ is the metric tree
length used by Notebook 18. All constants are uniform in
the earlier classical depth and periodic volume.

\begin{theorem}[Uniform classical-reference integration]
\label{ci-theorem}
Fix odd $L$, the gauge group, $c>0$ and $p>1/2$, and set
$R(g)=c(\log g^{-2})^p$. There are
$b,g_0,a_*,C,c_*>0$, independent of $M=L^k$ and compatible
periodic volume, for which
\begin{equation}\label{ci-normalized}
 Z_{M,g}(V)=\int d\mu_{A_M(V)^{-1}}(z)\,
       \chi_{R(g)}(z)\exp\left(\sum_Yv_{M,Y}(V,z)\right)
\end{equation}
has a holomorphic logarithm on the independent current
background product ball $\|a\|_\infty<b$, for $0<g<g_0$.
There is a convergent connected-support expansion
\begin{equation}\label{ci-result}
 \log Z_{M,g}(V)=\sum_HE_{M,g,H}(V|_H),\qquad
 \sup_i\sum_{H\ni i}e^{a_*|H|}\|E_{M,g,H}\|_b
       \le C\{gR(g)^3+e^{-c_*R(g)^2}\}.
\end{equation}
Its canonical difference from the pure Gaussian cutoff
expansion has norm at most $CgR(g)^3$.
\end{theorem}
\begin{proof}
Notebook 19 supplies the flat spectral constants and weighted
inverse uniformly in $M$. Choose the spatial exponents in
Notebook 18 smaller than both that inverse exponent and half
the available Steiner exponent for the support coefficients.
Choose a fixed $0<\theta\le1/4$, for example $1/(4D)$.
The arbitrary fixed cardinality reserve in the support
theorem pays $4+\log\mathcal K_\theta$ in its interaction norm.
The first bound in \eqref{ci-inputs} permits one
$M$-independent $b>0$ satisfying every local-precision
smallness condition of its tunable strictly local theorem.

The interaction parameter is at most $CgR^3$ by the second
bound in \eqref{ci-inputs}. The bad-cell parameter is at most
$Ce^{-c_*R^2}$ by the correlated real-baseline comparison.
The retained-bond rectangles are covered by the cell tail
argument of Notebook 20 and remain fixed under the exact
chart. A common $g_0$ makes both quarter budgets hold and
keeps the recovered whole field within its analytic radius.
Apply Notebook 18 and its periodic extension.
\end{proof}

At a physical background $V_*(W)$, the sum of the separated
linear terms is exactly zero and the gauge constant is zero.
Thus \eqref{ci-normalized} is the actual fluctuation factor
of the constrained classical-reference integral in that
chart. On the larger product domain it is the explicitly
centered analytic extension, as in Notebook 20.

The normalization factors are equally controlled. With
$B_M(V)=A_M(V)-A_{M,0}$ and $C_{M,0}=A_{M,0}^{-1}$, set
\begin{equation}\label{ci-det}
 \mathcal Q_M(V)=-\tfrac12\log\det(A_M(V)A_{M,0}^{-1})
   =\tfrac12\sum_{n\ge1}\frac{(-1)^n}{n}
                         \operatorname{Tr}[(C_{M,0}B_M(V))^n].
\end{equation}
The closed-chain proof of Notebook 20 uses only weighted
covariance row sums and the precision-label norm.
Equation \eqref{ci-inputs} supplies those uniformly in $M$,
even though the precision supports are no longer bounded
in size. Joining each label by a minimal tree and each
covariance link by a shortest path gives the same bound.
Consequently $\mathcal Q_M$ has strictly local connected
activities with anchored norm at most $Cb$. The separate
delta factor
$\log[\mathcal D(V)/\mathcal D(1)]=-\sum_c\operatorname{Tr}
\log\pi_c(V)$ has fixed-range anchored norm at most $Cb$.
Neither normalization term is silently counted as $O(g)$.

For real physical backgrounds define $I^{\rm sf}_{M,g}(V)$
to be the constrained current group integral with weight
$\exp[-g^{-2}(S_M(\log U/i)+\mathcal G(U))]$, restricted
to the local retained chart and its specified rectangle.
Using the notation $N_b,N_c,n,c_H,\mathcal D$ from
Notebook 20, its exact formula is
\begin{equation}\label{ci-exact}
 \begin{split}
 I^{\rm sf}_{M,g}(V)={}&
 c_H^{N_b-N_c}g^n\mathcal D(1)
       (2\pi)^{n/2}(\det A_{M,0})^{-1/2}
                       e^{-S_M(a)/g^2}\\
 &\times\exp\left\{
 \log\frac{\mathcal D(V)}{\mathcal D(1)}
       +\mathcal Q_M(V)+\sum_HE_{M,g,H}(V|_H)\right\}.
 \end{split}
\end{equation}
There is no unaccounted field-dependent normalization in
this defined integral. A gauge-insertion identity outside
the chart would additionally require $S_M$ itself to be
defined and gauge invariant on the chosen admissible
rooted-gauge-invariant domain. The small-potential construction
does not assert that extension. The exact chart formula above
does not require it.

\section{The next induction problem}
\label{ci-next}

The classical-reference integral is now controlled uniformly
over the number of earlier classical averaging steps.
This removes a separate localization premise from the
homogeneous small-field argument. Its support statement
initially records dependence on independent current potentials.
Notebook 22 supplies the further strict support expansion of
the actual stationary background $V_*(W)$ and transports the
generated activities into $W$, with explicit radius and
support-weight losses. Notebook 23 proves local gauge
invariance of the complete correction, including the
normalization factors, and its quadratic Ward identity.

Notebook 24 then supplies a conditional inherited-interaction
bound, and Notebook 26 meets its support-weight test by using
a persistent auxiliary-cell convention. The next problem is
stable analytic domains and control of the accumulated
interaction sizes, with extraction of vacuum and marginal
terms and estimates along the required coupling trajectory.
The bounded one-step estimates do not supply this recurrence.
Notebook 25's smaller representative gives one analytic
continuation tool for addressing the domain problem.
Large-field complements and mixed histories must be included
in the same induction. The continuum observable bounds,
symmetry restoration and reconstruction remain the final
existence obligations recorded in Notebook 01.

\begin{thebibliography}{9}
\bibitem{prior21}
\emph{Toward a four-dimensional Yang--Mills construction},
research Notebooks 12--20 in this workspace.
The principal inputs are the equal-exponent constrained
Schauder estimate, corrected nonlinear constraint, exact
Wilson force, nonlinear minimum, action-envelope identity,
flat retained covariance, overlap integration theorem and
exact physical source chart.
\bibitem{balaban10921}
T.~Ba\l aban,
\emph{Renormalization group approach to lattice gauge field
theories I}, Commun. Math. Phys. \textbf{109} (1987), 249--301,
\href{https://scispace.com/pdf/renormalization-group-approach-to-lattice-gauge-field-3ygai0vxt1.pdf}{original article copy}.
\bibitem{balaban11621}
T.~Ba\l aban,
\emph{Renormalization group approach to lattice gauge field
theories. II. Cluster expansions},
Commun. Math. Phys. \textbf{116} (1988), 1--22,
\href{https://scispace.com/pdf/renormalization-group-approach-to-lattice-gauge-field-4penm4k7hc.pdf}{original article copy}.
\bibitem{source21}
E.~Sussman and collaborators,
\href{https://github.com/ethanwsussman/Yang--Mills-LLM-experiment}{Yang--Mills-LLM-experiment},
deposited raw accounts at commit
\texttt{e392ef4501\allowbreak 3a581664ba\allowbreak b76b908671\allowbreak 70b30fe10e}.
\end{thebibliography}
\end{document}
